% =====================================================================
% BASELINE DE LITERATURA  ---  RESPONSABLE: INTEGRANTE 3
% Rubrica: 3 pts. Seccion propia para que se encuentre facil.
% La mitad de los puntos esta en la comparabilidad, no en el numero.
% =====================================================================
\section{Literature Baseline}
\label{sec:baseline-lit}

\subsubsection*{\textbf{Primary literature baseline}}
The best jurisdiction-matched published result we found is Silverman~\cite{ejeA_gupta2023}:
a hierarchical Bayesian model of running speed fit to 2,973 Hong Kong races
(36,006 observations), reaching 21.6\% accuracy in identifying the horse that
actually turns out fastest. We use this result as the primary literature reference
point for our evaluation, particularly for a per-race top-1 prediction task.
However, the value should not be interpreted as a strictly like-for-like
performance threshold because Silverman's prediction target is the horse with the
highest estimated running speed rather than the official race winner recorded in
our dataset. The 21.6\% result therefore serves as a jurisdiction-matched reference
point against which the difficulty of the prediction problem can be contextualized,
rather than as a directly equivalent benchmark.

We document the following comparability gaps rather than treating the number as a
like-for-like target, the main differences are as follow:
\begin{itemize}
  \item Race count and period: 2,973 races compared with our 6,349 races.
    The exact calendar period of Silverman's dataset is not stated in the source
    available, so we cannot determine whether it overlaps with, precedes, or
    follows our 1997--2005 period observation window.
  \item Target definition: Silverman's model predicts the horse with the
    highest estimated running speed, whereas our target is the race winner recorded
    in our dataset; these outcomes are related but not identical, since the horse
    estimated to be fastest may not necessarily finish first in a particular race.
  \item Technique: Silverman uses a hierarchical Bayesian model for
    running speed rather than a discriminative classifier or a learning-to-rank
    model. Consequently, differences in model family, target definition, and
    evaluation protocol prevent a direct apples-to-apples comparison without
    reproducing both approaches on the same dataset.
\end{itemize}

We also note, without using it as a baseline, that Silverman and
Suchard~\cite{ejeA_gupta2023} later extended this line of work toward a
discriminative conditional logit with frailty terms; it reports a computational
speedup (GPU vs. OpenMP by a factor of 2.68) for that model but no accuracy or AUC
figure, so it cannot be used quantitatively here (Section~\ref{sec:ejeB}).

\textbf{Secondary literature baseline)}
Chavan et al.~\cite{ejeC_chavan2024} provide the closest dataset-matched
reference identified in our review. The authors explicitly state that they use
the Hong Kong horse-racing dataset available on Kaggle, consisting of the
\texttt{races.csv} and \texttt{runs.csv} files, with 6,348 races and 79,447
horse records. Their preprocessing merges the two files using the race
identifier and encodes categorical variables such as venue and going before
model training. The study evaluates a Dummy Classifier, Logistic Regression,
Decision Tree, Random Forest, and Support Vector Machine using a random
80/20 train-test split. Random Forest achieves the highest reported precision
(62\%) and an accuracy of 92\%, with 52\% recall and a 51\% F1-score. The
authors also show that the Dummy Classifier achieves the same 92\% accuracy by
always predicting the majority \textit{not winning} class, illustrating why
accuracy alone is insufficient for this imbalanced prediction problem.

\subsubsection*{\textbf{Secondary analogs}} As additional reference
points are provided by studies conducted in racing environments outside Hong
Kong. Gupta and Singh~\cite{ejeA_gupta2023} evaluate Random Forest with SMOTE
on Turf Club of India data and report 93.1\% accuracy and 97.6\% ROC-AUC.
Davoodi and Khanteymoori~\cite{ejeA_davoodi2010} evaluate five neural-network
algorithms on 100 races from Aqueduct and report an average accuracy of 77\%.
Schumaker~\cite{ejeA_schumaker2018}, working with harness racing at Northfield
Park, reports 50.4\% wagering accuracy using support vector regression. These
results provide useful methodological context, but they are not directly
comparable with the present study because they differ in jurisdiction, racing
sub-type, dataset size, prediction target, and evaluation protocol. In
particular, wagering accuracy and winner-classification accuracy should not be
interpreted as equivalent measures. We therefore use these studies only as
contextual evidence rather than as quantitative performance targets.